%% kapa: exam=1 round=34 year=2023 period=2 form=A subject=appraisal_law subject_ko=감정평가관계법규 seq=31 number=31
%% source: Q-Net 감정평가사 기출문제 (cst00309) artlSeq=5211789 file=raw/1008/5211789/2023년도 제34회 감정평가사 1차시험 2교시.pdf
%% conversion: kordoc 4.24.51 -> tools/gen1.py (auto); review=proofread-sonnet; flags=none
\begin{question}{31}
공간정보의 구축 및 관리 등에 관한 법령상 지목에 관한 설명으로 옳은 것은?
\begin{choices}
  \item 토지가 임시적인 용도로 사용될 때에는 지목을 변경할 수 있다.
  \item 합병하려는 토지의 지목이 서로 다르더라도 소유자가 동일한 경우에는 토지의 합병을 신청할 수 있다.
  \item 자동차 정비공장 안에 설치된 급유시설 부지의 지목은 “주유소용지”로 한다.
  \item 고속도로의 휴게소 부지의 지목은 “도로”로 한다.
  \item 토지소유자는 지목변경을 할 토지가 있으면 그 사유가 발생한 날부터 30일 이내에 지적소관청에 지목변경을 신청하여야 한다.
\end{choices}
\end{question}
\answer{4}
